\documentclass[11pt]{article}

\usepackage[margin=1.05in]{geometry}
\usepackage[T1]{fontenc}
\usepackage{lmodern}
\usepackage{microtype}
\usepackage{amsmath,amssymb,amsthm}
\usepackage{booktabs}
\usepackage{array}
\usepackage{xcolor}
\usepackage[hidelinks]{hyperref}
\usepackage{enumitem}

\newtheorem{theorem}{Theorem}
\newtheorem{proposition}[theorem]{Proposition}
\newtheorem{lemma}[theorem]{Lemma}
\newtheorem{remark}[theorem]{Remark}
\newcommand{\C}{\mathbb C}
\newcommand{\A}{\mathbb A}
\newcommand{\Jac}{\operatorname{Jac}}
\newcommand{\Span}{\operatorname{span}}

\title{\textbf{Excluding the exact-\(\delta=4\) stratum}\\
\large in a quartic three-dimensional Keller leading family}
\author{Independent AI-assisted research program}
\date{First release: 26 July 2026 (UTC)}

\begin{document}
\maketitle

\begin{abstract}
Let \(F:\A^3_{\C}\to\A^3_{\C}\) be a polynomial map of total degree
four with nonzero constant Jacobian.  We consider the frozen
fixed-quadratic line-double-cover leading family
\[
H_4=(h(p,q)p^2,h(p,q)q^2,0),
\]
where \(h\) is a binary quadratic.  If \(R=(H_3)_3\), set
\[
\alpha=J(Q,R),\quad \beta=-J(P,R),\quad \gamma=J(P,Q),\quad
\delta=\deg\gcd(\alpha,\beta,\gamma).
\]
We prove that no Keller counterexample in this family has
\(\delta=4\).  A Hilbert--Burch and incidence calculation gives exactly
six normal-form orbits at \(\delta=4\).  Full lower-term determinant
descents exclude all six, with specialization-safe boundary charts.
The result is a structural exclusion inside one frozen quartic row.  It
does not prove a universal quartic exclusion and does not improve the
known total-degree floor four.
\end{abstract}

\noindent\textbf{Verification and review status.}
The proof and programs were produced with substantial language-model
assistance.  The work is not peer reviewed.  Exact symbolic checks certify
the encoded algebra, not the correctness of an unreviewed mathematical
model or worldwide priority.

\section{Scope and weighted determinant equations}

Write a degree-four map, after subtracting its constant term, as
\[
F=L(p,q,r)^t+H_2+H_3+H_4,
\]
where \(L\in\operatorname{GL}_3(\C)\) and \(H_i\) is homogeneous of
ordinary degree \(i\).  Introduce
\[
\Delta(w)=\det\left(L+wJH_2+w^2JH_3+w^3JH_4\right)
          =\sum_{j=0}^9 E_jw^j.
\]
The Keller condition is equivalent to
\[
E_0=\det L\ne0,\qquad E_j=0\quad(1\le j\le9).
\tag{1}
\]
We work only in the binary fixed-quadratic line-double-cover row
\[
H_4=(P,Q,0),\qquad P=hp^2,\qquad Q=hq^2,\qquad
h\in\C[p,q]_2,
\tag{2}
\]
and write \(R=(H_3)_3\in\C[p,q]_3\).  The top identity \(E_7=0\)
is controlled by the binary minors
\[
\alpha=J(Q,R),\qquad \beta=-J(P,R),\qquad \gamma=J(P,Q).
\tag{3}
\]
Here \(J\) denotes the two-variable Jacobian in \(p,q\).

\begin{theorem}\label{thm:main}
In the family \emph{(2)}, there is no quartic Keller counterexample for
which
\[
\delta:=\deg\gcd(\alpha,\beta,\gamma)=4.
\]
\end{theorem}

\begin{remark}
The theorem concerns one leading-shape stratum.  Other strata in the same
fixed-quadratic row remain open.  In particular, Theorem~\ref{thm:main}
does not imply that every dimension-three Keller counterexample has total
degree at least five.  The existing universal floor four is due to the
degree-three theorem known from Moh--Sathaye, for which Vistoli gave a short
conceptual proof \cite{Vistoli1999}.
\end{remark}

\section{Why exactly six families occur}

Put \(g=\gcd(\alpha,\beta,\gamma)\) and \(d=5-\delta\).  After dividing
by \(g\), the three generator degrees are
\[
(d,d,d+1).
\tag{4}
\]
The following bound makes the high-incidence enumeration finite.

\begin{lemma}[Hilbert--Burch deficit bound]\label{lem:hb}
Assume the reduced minor row has height two and its first two entries are
constant-linearly independent.  If \(e_1,e_2\) are its two minimal
syzygy degrees, then
\[
e_1+e_2=3d+1.
\]
The two gradient columns of \((P,Q,R)\) give syzygies of degree \(d+3\).
If \(k_i=d+3-e_i\), then
\[
0\le k_i\le2,\qquad k_1+k_2=\delta.
\tag{5}
\]
Consequently \(\delta\le4\); when \(\delta=4\), one has
\((k_1,k_2)=(2,2)\).
\end{lemma}

\begin{proof}
Hilbert--Burch applied to a height-two row of degrees \((d,d,d+1)\)
gives \(e_1+e_2=3d+1\).  Express the two gradient syzygies in a minimal
syzygy basis.  Their \(2\times2\) coefficient determinant is a nonzero
scalar multiple of \(g\): wedging the gradient columns gives
\((\alpha,\beta,\gamma)\), while wedging a minimal basis gives the reduced
row.  The coefficient degrees are \(k_i=d+3-e_i\).  They are nonnegative,
their sum is \(\deg g=\delta\), and the degree pattern in (4) bounds each
by two.  This proves (5) and the conclusions.
\end{proof}

Since
\[
\gamma=J(hp^2,hq^2)=8h^2pq,
\tag{6}
\]
the divisor \(g\) is selected from a finite list of exponent vectors on
the roots of \(h\), together with \(p=0\) and \(q=0\).  Divisibility of
\(J(Q,R)\) and \(J(P,R)\) gives a finite linear system for the four
coefficients of \(R\).  Solving those systems and quotienting by the
stabilizer of the double cover gives the following proposition.

\begin{proposition}[Exact-\(\delta=4\) incidence atlas]\label{prop:atlas}
Up to the frozen cover stabilizer, the independent exact-\(\delta=4\)
locus in \emph{(2)} is the disjoint union of the following six orbits.
\end{proposition}

\begin{center}
\small
\begin{tabular}{@{}>{\ttfamily}l p{0.55\textwidth} l@{}}
\toprule
\normalfont ID & \normalfont Normal form or modulus & \normalfont \(g\)\\
\midrule
D4-SF-21C  & \(h=XY,\ R=X^2Y,\ \kappa=-16/5\) & \(pX^2Y\)\\
D4-SF-20CC & squarefree \(h\), two fixed contacts, \(\kappa=16/5\)
            & \(pqX^2\)\\
D4-SF-11CC & squarefree \(h\), \(R=h(p+q)\), \(\kappa=16\)
            & \(pqh\)\\
D4-DN-3    & \(h=L^2,\ R=L^3\), \(L=p+q\) & \(L^4\)\\
D4-DN-2C   & \(h=L^2,\ R=L^2(p-2q)\), \(L=p+q\) & \(pL^3\)\\
D4-DN-1CC  & \(h=L^2,\ R=L(2p^2+pq+2q^2)\), \(L=p+q\)
            & \(pqL^2\)\\
\bottomrule
\end{tabular}
\end{center}

\begin{proof}
Factor \(h\) either as two distinct roots or as \(L^2\).  Equation (6)
bounds the exponent of every possible prime factor of \(g\).  For each
bounded exponent vector of total degree four, impose
\[
g\mid J(Q,R),\qquad g\mid J(P,R)
\tag{7}
\]
on a general binary cubic \(R\).  Empty systems are discarded; systems
whose actual gcd is larger or smaller are routed to the corresponding
exact stratum.  In the doubled-root chart, (7) leaves the three displayed
normal forms.  In the squarefree chart, the reciprocal and branch-swap
symmetries identify the surviving parameter values into the three
displayed orbits.  The exact enumeration checks all bounded exponent
vectors and all exit values before quotienting.

Two independent ledgers implement this calculation.  The primary ledger
has \(17+6+1\) high-incidence entries; a blinded reconstruction refines
the exact-\(\delta=3\) part and gives the canonical \(19+6+1\) ledger.
They agree on all six exact-\(\delta=4\) orbits.  The machine-readable
denominators, boundary routing, and strict reconciliation checks accompany
this paper.
\end{proof}

\section{The six lower-term exclusions}

For every orbit in Proposition~\ref{prop:atlas}, the calculation begins
with the complete \(E_7\) kernel, forms the full \(E_6\) system with every
lower coefficient present, and covers each rank or pivot boundary
separately.  The terminal mechanisms are summarized below.

\begin{center}
\small
\begin{tabular}{@{}>{\ttfamily}l p{0.69\textwidth}@{}}
\toprule
\normalfont Family & \normalfont Complete descent\\
\midrule
D4-SF-21C &
The \(E_6\) contact locus is a plane.  A rank-seven open chart and its
rank-six directions fail at \(E_5\); the rank-five origin loses all
nonbinary nonlinear coefficients and exits through a plane Keller map.\\
D4-SF-20CC &
The \(r^2\)-contact coordinates vanish and \(E_6\) leaves one contact
line.  Its nonzero chart fails at \(E_5\); the origin collapses at \(E_4\)
and exits through a plane Keller map.\\
D4-SF-11CC &
The complete contact plane splits into specialization-safe rank charts.
The nonzero charts fail in \(E_5,E_4\); the zero chart gives the same
unconditional plane exit.\\
D4-DN-3 &
The contact radical is two conjugate planes.  Their interiors fail at
\(E_5\), their punctured intersection forces \(\det L=0\), and the origin
collapses to the plane exit.\\
D4-DN-2C &
The contact radical is again two conjugate planes.  Their interiors fail
at \(E_5\).  On the punctured intersection, the full \(E_5\) residual
ideal splits into two branches; one forces \(\det L=0\), while the other
has a nonzero \(E_3\) residual.  The origin collapses to the plane exit.\\
D4-DN-1CC &
\(E_6\) leaves one contact line.  Its nonzero chart is inconsistent at
\(E_4\); at zero contact every nonlinear term is binary and the plane exit
applies.\\
\bottomrule
\end{tabular}
\end{center}

For completeness, we record the decisive identities in the last new
family.  On either transverse plane, with \(\eta^2=-2\), two
lower-variable-free \(E_5\) coefficients are
\[
\frac{(2k+3s)Q_1(k,s)}{162},\qquad
\frac{(2k+3s)Q_2(k,s)}{243},
\tag{8}
\]
and an exact Bezout identity shows that \(Q_1(1,t),Q_2(1,t)\) are
coprime.  Their common line is \(2k+3s=0\).  On its punctured part
\(k\ne0\), the full \(E_5\) residual ideal is
\[
\langle {\cal Q},{\cal A}{\cal B}\rangle.
\tag{9}
\]
The \(r^2\)-part of \(E_4\) forces \({\cal S}=0\), after which
\({\cal Q}={\cal D}^2\).  On \({\cal B}=0\), the remaining
compatibility is \({\cal A}{\cal F}_B=0\); the
\({\cal F}_B=0\) branch has \(\det L=0\), and
\({\cal A}=0\) is the other branch.  There,
\[
\det L=\frac{{\cal F}_A{\cal H}_A}{1152}.
\tag{10}
\]
If \({\cal F}_A\ne0\), two \(E_3\) equations have pivot
\[
\frac{k^2{\cal F}_A^2}{144},
\]
while each remaining equation reduces to
\[
\frac{k{\cal F}_A^2}{288}\ne0.
\tag{11}
\]
At the origin, a fresh constant \(E_6\) pivot gives
\[
[p^3r]E_4=-3b_{qr}^2,\qquad
[q^3r]E_4=\frac23(3b_{qr}+2\ell_8)^2.
\tag{12}
\]
Thus all six nonbinary quadratic coefficients vanish.

Whenever a zero-contact chart leaves only binary nonlinear terms, a
linear target normalization separates a two-variable Keller map of
degree at most four and a triangular \(r\)-coordinate.  Moh's
unconditional plane theorem for degree below \(100\)
\cite{BiggersMohFried1983} makes that plane map an automorphism.  This is
a proved bounded-degree theorem, not an assumption of the open plane
Jacobian conjecture.

\begin{proof}[Proof of Theorem~\ref{thm:main}]
Proposition~\ref{prop:atlas} is exhaustive and disjoint.  The complete
lower-term descents summarized above exclude each of its six terminal
orbits.  Therefore no exact-\(\delta=4\) leading point in (2) can belong
to a quartic Keller counterexample.
\end{proof}

\section{Exact certificates and limitations}

The aggregate verifier performs the following checks.
\begin{enumerate}[leftmargin=2em]
\item It compares a six-entry proof manifest with the canonical blinded
      \(19+6+1\) denominator.
\item It deletes one family and requires the bridge verification to fail.
\item It reruns the primary enumeration, blinded denominator,
      reconciliation, and frozen hashes.
\item It obtains the execution plan from the validated manifest and runs
      six distinct fail-closed family wrappers.
\item Each family wrapper requires exact terminal markers from symbolic
      algebra and an independently implemented replay, and contains
      deliberate coefficient mutations that must fail.
\item A separate hostile assembly wrapper reruns the complete aggregate,
      verifies that the containing global row remains open, and rejects
      mutations of the canonical ID, alias, proof path, terminal marker,
      denominator source, theorem scope, and advertised counts.
\end{enumerate}
The terminal assembly marker is
\[
\texttt{EXACT\_DELTA4\_SIX\_FAMILY\_EXCLUSION\_STRICT\_PASS}.
\]
The hostile assembly terminates with
\[
\texttt{EXACT\_DELTA4\_HOSTILE\_UMBRELLA\_AUDIT\_STRICT\_PASS}.
\]

These checks are reproducible evidence for the stated stratum theorem.
They do not constitute peer review.  In particular, they do not close the
remaining exact-\(\delta=3\), lower-incidence, or power-fibre families,
the parent fixed-quadratic row, or the global quartic denominator.

\begin{thebibliography}{9}

\bibitem{BiggersMohFried1983}
R. Biggers, T. T. Moh, and M. Fried,
\emph{On the Jacobian conjecture and the configurations of roots},
J. Reine Angew. Math. \textbf{340} (1983), 140--213.

\bibitem{Vistoli1999}
A. Vistoli,
\emph{The Jacobian conjecture in dimension 3 and degree 3},
J. Pure Appl. Algebra \textbf{142} (1999), 79--89.

\end{thebibliography}

\end{document}
